Методы пассивной балансировки позволяют выравнивать напряжение только при полной зарядке батареи
и в простое ячейки не балансируются. Активные методы напротив не имеют этого недостатка.
Ток балансировки подбирается для конкретной сборки исходя из емкости одной ячейки.
Пасивный метод прост в реализации и не требует сложного програмного управления. 
Так же многогие крупные компании используют данный метод балансировки в своих продуктах для массового производства.


Итак в качестве практического упражнения была выбрана пассивная \acs{bms}.
Необходимо балансировать 8s батарею. Балансировочный ток 200 mA. 
Плата должна отслеживать напряжение и температуру каждой ячейки.
Данные отправляются в CAN сеть.
